\subsection{CAYLEY ALGEBRAS}

DEFINITION 1. A Cayley algebra over \(\mathbf{A}\) is an ordered pair \((\mathbf{E}, s)\) , where \(\mathbf{E}\) is an algebra over \(\mathbf{A}\) with a unit element \(e\) and \(s\) is an antiautomorphism of \(\mathbf{E}\) such that

\[
u + s (u) \in \mathrm{A} e \quad \text { and } \quad u. s (u) \in \mathrm{A} e
\]

for all \(u \in \mathbf{E}\) .

s is called the conjugation of the Cayley algebra (E, s) and \(s(u)\) the conjugate of u. The condition \(u + s(u) \in Ae\) implies that u and \(s(u)\) are permutable. We write

\begin{enumerate}
\def\labelenumi{(\arabic{enumi})}
\setcounter{enumi}{8}
\tightlist
\item
\end{enumerate}

\[
\mathrm{T} (u) = u + s (u)\tag{10}
\]

\[
\mathrm{N} (u) = u. s (u) = s (u). u
\]

and these elements of the subalgebra Ae are called respectively the Cayley trace and norm of u.

The ordered pair consisting of a quadratic algebra E and its conjugation s (which is an antiautomorphism since E is commutative) (no. 3) is a Cayley algebra.

Let \((\mathbf{E}, s)\) be a Cayley algebra; as \(s(e) = e\) , \(s(u + s(u)) = u + s(u)\) , in other words \(s(u) + s^2(u) = u + s(u)\) or also

\[
s ^ {2} (u) = u\tag{11}
\]

so that \(s^{2}\) is the identity mapping of E. It follows that

\[
\mathrm{T} (s (u)) = \mathrm{T} (u), \quad \mathrm{N} (s (u)) = \mathrm{N} (u).\tag{12}
\]

Finally, the relation \((u - u)(u - s(u)) = 0\) gives

\[
u ^ {2} - \mathrm{T} (u). u + \mathrm{N} (u) = 0\tag{13}
\]

for all \(u \in \mathbf{E}\) .

PROPOSITION 4. Let E be an A-algebra and s and \(s'\) antiautomorphisms of E such that \((E, s)\) and \((E, s')\) are Cayley algebras. If E admits a basis containing the unit element e, then \(s' = s\) .

Clearly \(s'(u) = s(u) = u\) for all \(u \in Ae\) . If T, N (resp. T', N') are the trace and norm functions for (E, s) (resp. (E, s')), it follows from (13) that

\[
(\mathrm{T} (u) - \mathrm{T} ^ {\prime} (u)). u - (\mathrm{N} (u) - \mathrm{N} ^ {\prime} (u)) = 0.
\]

Let B be a basis of E containing e and u an element of B distinct from e; then \(\mathrm{T}(u) - \mathrm{T}'(u) = 0\) , whence \(s(u) = s'(u)\) . As s and \(s'\) coincide on B, they are equal.

In what follows, we shall write \(\bar{u} = s(u)\) , so that

\[
\left\{ \begin{array}{l l} u + \bar {u} = \mathrm{T} (u), & u \bar {u} = \bar {u} u = \mathrm{N} (u), \quad \bar {\bar {u}} = u \\ \overline {{u + v}} = \bar {u} + \bar {v}, & \overline {{\alpha u}} = \alpha \bar {u}, \quad \overline {{u v}} = \bar {v}. \bar {u} \end{array} \right.\tag{14}
\]

for \(u, v\) in \(\mathbf{E}, \alpha \in \mathbf{A}\) ; moreover

\[
\mathbf {T} (e) = 2 e, \quad \mathbf {N} (e) = e.\tag{15}
\]

From the formula

\[
\mathrm{T} (u v) = u v + \overline {{{{u v}}}} = u v + \bar {v}. \bar {u} = u v + (\mathrm{T} (v) - v) (\mathrm{T} (u) - u),
\]

we deduce that

\[
u v + v u = \mathrm{T} (u) v + \mathrm{T} (v) u + (\mathrm{T} (u v) - \mathrm{T} (u) \mathrm{T} (v))\tag{16}
\]

whence, exchanging \(u\) and \(v\) ,

\[
\mathbf {T} (v u) = \mathbf {T} (u v).\tag{17}
\]

On the other hand, \(\mathrm{N}(u+v)=(u+v)(\bar{u}+\bar{v})=\mathrm{N}(u)+\mathrm{N}(v)+\mathrm{T}(u\bar{v})\) , whence

\[
\mathrm{T} (v \bar {u}) = \mathrm{T} (u \bar {v}) = \mathrm{N} (u + v) - \mathrm{N} (u) - \mathrm{N} (v).\tag{18}
\]

Now, (16) applied with u replaced by \(\bar{u}\) gives

\[
\mathrm{T} (u \bar {v}) = \mathrm{T} (u) \mathrm{T} (v) + \bar {u} v + v \bar {u} - \mathrm{T} (u) v - \mathrm{T} (v) \bar {u} = \mathrm{T} (u) \mathrm{T} (v) - u v - \bar {v}. \bar {u};
\]

whence

\[
\mathrm{T} (v \bar {u}) = \mathrm{T} (u \bar {v}) = \mathrm{N} (u + v) - \mathrm{N} (u) - \mathrm{N} (v) = \mathrm{T} (u) \mathrm{T} (v) - \mathrm{T} (u v).\tag{19}
\]

Finally, clearly for all \(\alpha \in \mathbf{A}\) ,

\[
\mathrm{N} (\alpha u) = \alpha^ {2} \mathrm{N} (u);\tag{20}
\]

in particular \(\mathbf{N}(2u) = 4\mathbf{N}(u)\) , so that formula (19) gives

\[
(\mathrm{T} (u)) ^ {2} - \mathrm{T} (u ^ {2}) = 2 \mathrm{N} (u).\tag{21}
\]

Clearly T is a linear form on the (Ae)-module E. As \((u, v) \mapsto \mathrm{T}(v\bar{u})\) is a bilinear form on this module, it follows from (18) and (20) that N is a quadratic form (cf.~IX, § 3, no. 4).

